\documentclass{article}
\usepackage[french]{babel}
\title{Un portfolio généré par l'IA, relu et livré proprement}
\author{Charles Lepaire}
\date{Octobre 2026}

\begin{document}
\maketitle

\begin{abstract}
Ce portfolio montre qu'un docteur en informatique graphique sait faire générer du code par l'IA dans
une quinzaine de langages, puis le relire, le tester et le livrer. Le fil rouge est celui de ma thèse :
les maillages 3D et leur topologie. Ce document est lui-même rendu par l'un des projets, une
bibliothèque \LaTeX{} écrite en TypeScript.
\end{abstract}

\tableofcontents

\section{Objectif}\label{sec:objectif}
Un recruteur voit souvent une liste de technologies sans preuve. Ici, chaque projet arrive avec ses
tests unitaires et d'intégration, une intégration continue verte sur Linux, Windows et macOS quand il
est compilé, une démonstration en ligne et deux fichiers de suivi : \texttt{DECISIONS.md} pour les
choix d'architecture, \texttt{REVIEW.md} pour la relecture humaine du code généré.

\section{Méthode}\label{sec:methode}
Le travail suit Scrum, par sprints de deux semaines. À chaque étape :
\begin{enumerate}
\item une spécification et un plan d'architecture, relus ;
\item les tests d'abord, puis la génération par petites étapes ;
\item la vérification qu'un test échoue quand on casse le code\footnote{Par exemple, sans le test
d'occurrence de l'inférence de types du mini-langage, les tests unitaires tombent.} ;
\item la relecture de chaque différence, et chaque défaut noté avec sa correction.
\end{enumerate}
Les défauts trouvés ainsi sont réels : des entiers de 32 bits sous js\_of\_ocaml, des exposants à
trois chiffres sous Windows, une pile JavaScript trop petite sous macOS.

\section{Le fil rouge : les maillages}\label{sec:maillages}
Les projets forment une chaîne. Une bibliothèque C lit les fichiers OBJ, PLY et STL ; une bibliothèque
C++ en construit la structure demi-arête et en calcule les invariants. Le premier est la
caractéristique d'Euler d'une surface à $V$ sommets, $E$ arêtes et $F$ faces :
\begin{equation}
\chi = V - E + F. \label{eq:euler}
\end{equation}
Pour une surface orientable à $c$ composantes et $b$ boucles de bord, elle donne le genre $g$, le
nombre de « trous » :
\begin{equation}
\chi = 2c - 2g - b. \label{eq:genre}
\end{equation}
La courbure de Gauss discrète en un sommet intérieur est le défaut angulaire
$K_v = 2\pi - \sum_i \theta_i$, et le théorème de Gauss-Bonnet relie la géométrie à la topologie :
\begin{equation}
\sum_v K_v = 2\pi\chi. \label{eq:gauss-bonnet}
\end{equation}
Les tests vérifient \eqref{eq:gauss-bonnet} à $10^{-9}$ près sur des tores déformés au hasard, et le
mini-langage Maille calcule \eqref{eq:genre} sur des scènes : \texttt{genus (union t (union t s))}
vaut 2 pour deux tores et une sphère. Ces notions viennent de ma thèse~\cite{these}, sur les
cartes généralisées et l'arbre vasculaire cérébral~\cite{wscg}.

\section{Les projets}\label{sec:projets}
\begin{tabular}{lll}
\hline
\textbf{Projet} & \textbf{Langages} & \textbf{Ce qui est prouvé} \\
\hline
Bibliothèque de maillages & C, WebAssembly & fuzzing, Valgrind, 3 systèmes \\
Topologie 3D & C++, three.js & Gauss-Bonnet, viewer en ligne \\
API & Python, FastAPI & Docker, invariants par HTTP \\
Apprentissage & PyTorch, MLflow & seuil de précision bloquant \\
Carrefour & Ada, SPARK & absence d'erreur prouvée \\
Base de mesures & SQL, PostgreSQL & pgTAP, plan d'exécution \\
Mini-langage & C, Bison, OCaml & types inférés \\
Cet éditeur & TypeScript, KaTeX & diagnostics situés \\
\hline
\end{tabular}

\noindent La liste complète, avec l'état de chaque projet, est sur la page Projets. La section
\ref{sec:methode} décrit comment chacun est relu.

\section{Ce que montre ce document}
Il est analysé et rendu dans votre navigateur. Modifiez une formule ou supprimez une accolade :
l'aperçu se met à jour et les diagnostics indiquent la ligne et la colonne du problème.

\begin{thebibliography}{9}
\bibitem{these} C. Lepaire. \emph{Extraction par analyse topologique sur les G-cartes de
caractéristiques pour les systèmes d'apprentissage supervisé appliqués à l'arbre vasculaire
cérébral}. Thèse de doctorat, Université de Poitiers, 2025.
\url{https://theses.fr/s342053}
\bibitem{wscg} C. Lepaire, H. Belhaouari, R. Pascual, P. Meseure. Exploitation of local adjacencies
for parallel construction of a Reeb graph variant: cerebral vascular tree case. \emph{Journal of
WSCG}, 2024. \url{https://doi.org/10.24132/JWSCG.2024.9}
\end{thebibliography}

\end{document}
